\documentclass[12pt,a4paper]{article}
\usepackage[utf8]{inputenc}
\usepackage[margin=1in]{geometry}
\usepackage{helvet}
\renewcommand{\familydefault}{\sfdefault}
\usepackage{enumitem}
\usepackage{hyperref}
\usepackage[dvipsnames]{xcolor}

\hypersetup{
    colorlinks=false,
    pdfborder={0 0 0}
}

\setlength{\parindent}{0pt}
\setlength{\parskip}{6pt}

\begin{document}

\section*{POIN PERBAIKAN IN-SCOPE}

\subsection*{A. PERFORMA \& INFRASTRUKTUR}

\textbf{1. JavaScript terlalu banyak (53 scripts)}

Terkait: Sistem Web Aplikasi Hal 4 Poin 2 (Cache gambar, CSS, JS, file statis)

\textbf{2. Network payload 2,63 MB (target: 1--1,5 MB)}

Terkait:
\begin{itemize}[topsep=0pt, partopsep=0pt, itemsep=2pt]
    \item Sistem Web Aplikasi Hal 4 Poin 2 (Integrasi CDN, image optimizer, caching)
    \item Sistem Web Aplikasi Hal 5 Modul 5 (Upload gambar auto resize \& compress max 1 MB)
    \item Kebutuhan User Hal 2 (Mobile Friendly, Upload foto Auto size 1200×800 Max 1 MB)
\end{itemize}

\textbf{3. Render-blocking (3 resources)}

Terkait: Sistem Web Aplikasi Hal 4 Poin 2 (Cache untuk Google Lighthouse > 95)

\textbf{4. Server response 546 ms (target: <300 ms)}

Terkait:
\begin{itemize}[topsep=0pt, partopsep=0pt, itemsep=2pt]
    \item Sistem Web Aplikasi Hal 4 Poin 2 (Next.js SSR/SSG, CDN)
    \item Sistem Web Aplikasi Hal 5 Modul 5 (Pre-signed URL Object Storage)
\end{itemize}

\textbf{5. Cache belum optimal}

Terkait: Sistem Web Aplikasi Hal 4 Poin 2 (Integrasi CDN, image optimizer, caching)

\subsection*{B. TECHNICAL SEO}

\textbf{6. Canonical mismatch (PRIORITAS TINGGI)}

Terkait:
\begin{itemize}[topsep=0pt, partopsep=0pt, itemsep=2pt]
    \item Sistem Web Aplikasi Hal 6 Modul 8 (Meta tag \& sitemap otomatis)
    \item Kebutuhan User Hal 3 (Situs website mudah terindeks Google)
\end{itemize}

\textbf{7. 8 gambar tanpa ALT}

Terkait: Sistem Web Aplikasi Hal 6 Modul 8 (Meta tag \& sitemap otomatis untuk SEO)

\textbf{8. 10 gambar tanpa width/height}

Terkait:
\begin{itemize}[topsep=0pt, partopsep=0pt, itemsep=2pt]
    \item Sistem Web Aplikasi Hal 5 Modul 5 (Upload gambar auto resize 1200×800 px)
    \item Kebutuhan User Hal 2 (Upload foto Auto size 1200×800 Megapixel)
\end{itemize}

\textbf{9. Error 503 di media.arasvara.id / configuration.arasvara.id}

Terkait:
\begin{itemize}[topsep=0pt, partopsep=0pt, itemsep=2pt]
    \item Sistem Web Aplikasi Hal 4 Poin 2 (Integrasi Object Storage dan subdomain)
    \item Sistem Web Aplikasi Hal 5 Modul 5 (Integrasi Object Storage Cloud)
\end{itemize}

\textbf{10. Heading structure: H3 missing}

Terkait: Sistem Web Aplikasi Hal 4 Poin 1 (Wireframe halaman utama, artikel, dashboard, CMS)

\textbf{11. NewsArticle schema belum lengkap/optimal}

Terkait: Sistem Web Aplikasi Hal 6 Modul 8 (Artikel terindeks optimal di mesin pencari)

\subsection*{C. CMS \& SECURITY}

\textbf{12. Kurangi akun Super Admin}

Terkait:
\begin{itemize}[topsep=0pt, partopsep=0pt, itemsep=2pt]
    \item Sistem Web Aplikasi Hal 5 Modul 4 (Pembatasan hak akses per level role)
    \item Kebutuhan User Hal 3 (CMS kategori Admin IT, Editor, Reporter + pembatasan level)
\end{itemize}

\textbf{13. Aria-label button ikon di header}

Terkait: Sistem Web Aplikasi Hal 4 Poin 1 (Tata letak dan elemen branding)

\textbf{14. Responsivitas tabel di mobile}

Terkait:
\begin{itemize}[topsep=0pt, partopsep=0pt, itemsep=2pt]
    \item Sistem Web Aplikasi Hal 6 Modul 6 (Tampilan cepat, ringan, responsif)
    \item Kebutuhan User Hal 2 (Mobile Friendly)
\end{itemize}

\textbf{15. Character counter di meta description field}

Terkait:
\begin{itemize}[topsep=0pt, partopsep=0pt, itemsep=2pt]
    \item Sistem Web Aplikasi Hal 5 Modul 2 (CRUD berita, preview, auto-save)
    \item Sistem Web Aplikasi Hal 6 Modul 8 (SEO \& Analitik)
    \item Kebutuhan User Hal 2 (CMS Simple dan modern)
\end{itemize}

\subsection*{D. UI/UX (MINOR)}

\textbf{16. Kontras warna link footer rendah}

Terkait: Sistem Web Aplikasi Hal 4 Poin 1 (Penyesuaian warna, tata letak, font, branding)

\textbf{17. Tombol pencarian mengambang menutupi konten mobile}

Terkait:
\begin{itemize}[topsep=0pt, partopsep=0pt, itemsep=2pt]
    \item Sistem Web Aplikasi Hal 4 Poin 1 (Rancangan visual \& navigasi)
    \item Sistem Web Aplikasi Hal 6 Modul 6 (Mobile friendly/responsif)
    \item Kebutuhan User Hal 2 (Mobile Friendly)
\end{itemize}

\textbf{18. Heading kategori: pisahkan ikon panah dari teks}

Terkait: Sistem Web Aplikasi Hal 4 Poin 1 (Wireframe dan prototipe interaktif di Figma)

\subsection*{E. MONITORING \& INVESTIGATION (PARTIAL)}

\textbf{19. Backlink domain tidak dikenal}

Developer scope:
\begin{itemize}[topsep=0pt, partopsep=0pt, itemsep=2pt]
    \item Setup monitoring/audit tool untuk track backlink profile
    \item Generate disavow file jika dibutuhkan
\end{itemize}

Tim Arasvara scope:
\begin{itemize}[topsep=0pt, partopsep=0pt, itemsep=2pt]
    \item Review backlink list dan tentukan mana yang spam/suspicious
    \item Submit disavow list ke Google Search Console
\end{itemize}

\textbf{20. Indexing Google lambat}

Developer scope:
\begin{itemize}[topsep=0pt, partopsep=0pt, itemsep=2pt]
    \item Verify sitemap submission \& robots.txt
    \item Check crawl budget issues (server response, redirect chains)
    \item Implement Google Indexing API jika diperlukan
    \item Fix technical SEO barriers (canonical, speed, mobile-friendly)
\end{itemize}

Tim Arasvara scope:
\begin{itemize}[topsep=0pt, partopsep=0pt, itemsep=2pt]
    \item Strengthen internal linking antar artikel
    \item Improve content quality \& E-E-A-T signals
    \item Add primary source outbound links
    \item Regular content publishing schedule
\end{itemize}

\vspace{12pt}
\hrule
\vspace{12pt}

TOTAL IN-SCOPE: 20 poin (18 full developer + 2 partial)

\vspace{12pt}
\hrule
\vspace{12pt}

\section*{OUT OF SCOPE (Perlu Klarifikasi Client)}

\textbf{21. Pagination multi-halaman artikel}
\begin{itemize}[topsep=0pt, partopsep=0pt, itemsep=2pt]
    \item Permintaan client untuk otomatis, tidak diubah
\end{itemize}

\textbf{22. 2FA}
\begin{itemize}[topsep=0pt, partopsep=0pt, itemsep=2pt]
    \item Tidak ada pembahasan kontrak
\end{itemize}

\textbf{23. Revisi UI major lebih dari 1x}
\begin{itemize}[topsep=0pt, partopsep=0pt, itemsep=2pt]
    \item Jika perlu re-layout author page/about us, cuma 1 kali, no major revision lagi
\end{itemize}

\vspace{12pt}
\hrule
\vspace{12pt}

\section*{TIM ARASVARA (Content \& Editorial)}

\textbf{24. Internal linking lemah (14 links)}
\begin{itemize}[topsep=0pt, partopsep=0pt, itemsep=2pt]
    \item Editor wajib tambahkan internal link relevan antar artikel
    \item Contoh: artikel harga emas → link ke artikel emas hari sebelumnya, cara beli emas, investasi pemula
\end{itemize}

\textbf{25. 0 external links (MERAH)}
\begin{itemize}[topsep=0pt, partopsep=0pt, itemsep=2pt]
    \item Wajib link ke sumber primer: Antam, BI, BPS, kementerian, perusahaan, dokumen resmi
    \item Jangan cuma tulis ulang media lain
\end{itemize}

\textbf{26. Homepage text 224 kata (kuning)}
\begin{itemize}[topsep=0pt, partopsep=0pt, itemsep=2pt]
    \item Tambah deskripsi konten/intro jika perlu
\end{itemize}

\textbf{27. E-E-A-T belum kuat}
\begin{itemize}[topsep=0pt, partopsep=0pt, itemsep=2pt]
    \item Buat halaman author per penulis
    \item Format: Nama, jabatan, bio singkat, link ``Lihat semua artikel [Nama]''
    \item Contoh: ``Nurul Kharimah adalah reporter Arasvara yang meliput lifestyle, entertainment dan isu generasi muda''
\end{itemize}

\textbf{28. Author markup di artikel}
\begin{itemize}[topsep=0pt, partopsep=0pt, itemsep=2pt]
    \item Hubungkan penulis artikel dengan URL halaman profil mereka
\end{itemize}

\textbf{29. Typo di judul artikel yang sudah tayang}
\begin{itemize}[topsep=0pt, partopsep=0pt, itemsep=2pt]
    \item QC editorial sebelum publish
\end{itemize}

\textbf{30. NAP (Name, Address, Phone) tidak ditemukan di page biasa}
\begin{itemize}[topsep=0pt, partopsep=0pt, itemsep=2pt]
    \item Cantumkan di footer/contact selain About Us
\end{itemize}

\end{document}
